\maketitle
\tableofcontents

% ============================================================
\chapter*{Doc.~308 -- $\Lambda$ as a Reading Artefact}
\addcontentsline{toc}{chapter}{Doc.~308 -- $\Lambda$ as a Reading Artefact}
% ============================================================

% ============================================================
\section*{Preface and Context}
% ============================================================

This document has two cores that belong together.

\medskip
\noindent\textbf{Core 1 — The Two-Halves Synthesis (Stage B):}
Einstein 1911 calculated the deflection of light at the Sun from
time dilation alone: $\alpha = 0.875''$ — half the measured value.
Einstein 1915 (GR) yielded $\alpha = 1.750''$, because space
\emph{and} time curvature each contribute one half.

In FFGFT, $\tilde T \cdot m = 1$ (time curvature) and fractal path
elongation (Doc.~182, space curvature) are \emph{not} alternatives
but \textbf{complementary halves of the same deflection} — each
exactly $0.875''$, together exactly $1.750''$. This is not a fit but
a consequence of K3 area invariance ($\gamma_\text{PPN} = 1$ exactly,
Cassini passed). FFGFT \emph{explains} why Einstein 1911 was missing
half: he had only the time rate, not the space half.

\medskip
\noindent\textbf{Core 2 — Ontological Status of the Dark Sector:}
If redshift admits two readings — expansion ($a(t)$) or static
geometric path stretching ($1+z = e^{\xi x}$, Doc.~280) — and the
classical data (SNe~Ia, BAO, CMB) are degenerate between both
(Doc.~267), then quantities that exist \emph{only within one reading}
have the status of bookkeeping quantities without reading-independent
reference.
Dark energy satisfies all three artefact conditions. Dark matter does
not — it has reading-independent evidence (rotation curves, Bullet
Cluster) and is an open geometric problem for FFGFT.

\medskip
\noindent\textbf{Connection of both cores:} The space half of the
Einstein deflection ($\tilde T \cdot m = 1$ + K3 area) is structurally
the same as the potential that flattens rotation curves at accelerations
$a < a_0 = cH_0/(2\pi)$ — $a_0$ comes from the same $\xi^{10}$ chain
(Stage~A). Dark matter and dark energy are therefore not equally open:
DM has geometric candidates in the corpus; DE has no reading-independent
existence.

This document is not part of the A-series. It references
Doc.~267, 280, 182, 279, 275 and Doc.~190 (P18--P20, P33, P35, P44, K6).

\medskip
\noindent\textbf{[POSIT] Thesis:} Dark energy is a reading artefact —
computationally present, without reading-independent consequence.
Dark matter is \emph{not} — it has reading-independent evidence and
geometric candidates (P44 A–C computed, $\mu$ derivation from T$^4$
open).

% ============================================================
\section*{Conventions}
% ============================================================

\noindent\textbf{Legend.} Every substantive claim carries a layer
label in square brackets:

\medskip
\begin{tabular}{@{}lp{11cm}}
\textbf{[K]}      & \emph{Core derivation} — derived mathematically
                    or numerically from the foundational posits;
                    internally consistent and reproducible. \\[4pt]
\textbf{[S]}      & \emph{Structural claim} — algebraically or
                    geometrically motivated, but not yet fully proved. \\[4pt]
\textbf{[POSIT]}  & \emph{Posit} — declared assumption or definition
                    that is not further derived; must be accepted as
                    such or justified separately. \\
\end{tabular}

\medskip
\noindent\textbf{P-numbers and K-numbers} come from the FFGFT
correction register Doc.~190 (project register). \textbf{P}$n$
denotes an open programme point (a gap still to be closed or an
ongoing computation); \textbf{K}$n$ denotes a completed correction
or anchored finding. Relevant entries in this document:

\medskip
\begin{tabular}{@{}lp{10.5cm}}
\textbf{P44}      & MOND/rotation curves/Bullet Cluster: three-stage
                    FE programme A/B/C for a geometric explanation of
                    rotation curves and gravitational lensing without
                    particle DM. \\[4pt]
\textbf{P33}      & Magnitude of $\xi = 4/30000$ ($5^4$ scale):
                    declared intuitively/empirically motivated, no
                    forward derivation. \\[4pt]
\textbf{P35}      & Bulk exponent $k^*$: $O(1)$ prefactor still open. \\[4pt]
\textbf{K6}       & $H_0$ as $\xi^{10}$ power: anchored,
                    $a_0$ coincidence derived. \\[4pt]
\textbf{P18--P20} & Earlier open points in the cosmological sector
                    (degeneracy, redshift, static reading). \\
\end{tabular}

% ============================================================
\chapter{Definition: Reading Artefact}
% ============================================================

\textbf{[POSIT]} A quantity is called a \emph{reading artefact} here
when it simultaneously satisfies three conditions:

\begin{enumerate}
  \item It is determined exclusively from data that are degenerate
        between the competing readings (Doc.~267).
  \item It possesses no independent anchor outside this dataset —
        no laboratory measurement, no local dynamics, no obligation
        in a foreign sector.
  \item When the reading changes, it does not disappear as
        \emph{refuted} but as \emph{objectless}: there is nothing
        it referred to.
\end{enumerate}

Historical parallel: epicycles. They computed correctly, within their
reading, and yet referred to nothing. The transition to Kepler did not
falsify them but rendered them objectless. A reading artefact is not
a computational error — it is a quantity without referent.

% ============================================================
\chapter{Dark Energy: All Three Conditions Satisfied}
% ============================================================

\textbf{[K] Condition 1.} All evidence for $\Lambda$ is cosmological:
SNe~Ia luminosities, BAO angular scales, CMB acoustic peaks —
exactly the three pillars whose degeneracy between the expansion and
static readings Doc.~267 demonstrates in detail. No source of evidence
for $\Lambda$ exists outside this dataset.

\medskip
\textbf{[K] Condition 2.} $\Lambda$ has no independent anchor —
and more sharply: the only theoretical candidate has failed. The
QFT vacuum energy density misses the fit value by $\sim 10^{122}$
(the fine-tuning problem; see Doc.~279, outlook). $\Lambda$CDM thus
introduces a quantity that (a) is not reading-independently forced by
the degenerate data and (b) is missed by its own underlying theory by
122 orders of magnitude. $\Lambda$ carries no obligation in any other
sector of physics: it works nowhere except its own fit.

\medskip
\textbf{[K] Condition 3.} In the static reading (Doc.~280) there is
no expansion to be accelerated. The apparent acceleration is part of
the degeneracy structure (Doc.~267, §\,accelerated dilatation in both
models). When the reading changes, $\Lambda$ is not refuted — it loses
its object.

\medskip
\noindent\textbf{[K] Conclusion:} Dark energy is a reading artefact
in the sense of Ch.~1. Formulation with care: \emph{not} ``there is no
dark energy'' (that would be a statement within the expansion reading),
but: ``dark energy is a reading-internal bookkeeping quantity without
reading-independent existence evidence.''

% ============================================================
\chapter{Dark Matter: Condition 1 Violated — Not an Artefact}
% ============================================================

\textbf{[K]} Dark matter does \emph{not} satisfy Condition~1.
Its evidence lies outside the redshift degeneracy:

\begin{itemize}
  \item \textbf{Rotation curves:} local dynamics of individual galaxies
        — independent of the $z$ reading.
  \item \textbf{Gravitational lensing:} geometry of light propagation
        at individual masses — independent of the $z$ reading.
  \item \textbf{Bullet Cluster:} spatial separation of lensing mass
        and X-ray gas in a single object — independent of the $z$ reading.
\end{itemize}

These findings do not disappear when the reading changes. They require
an explanation in \emph{every} reading.

\medskip
\noindent\textbf{[POSIT]} For FFGFT, dark matter is therefore not an
artefact but an \textbf{open geometric problem}: the candidate is
fractal path elongation (Doc.~182) as a source of additional effective
attraction without a particle substrate — status: reduction layer, no
core derivation. P44 Stages~A–C are computed: $a_0$ anchor (Stage~A),
factor-2 gate and selection computation (Stage~B), Bullet Cluster and
rotation curves (Stage~C). Remaining open: forward proof of K3 and
form of $\mu(x)$ from T$^4$ geometry (P44 in Doc.~190).

\medskip
\noindent\textbf{Caution regarding one's own framework:} The formulation
``the dark sector is invented'' would be sweeping and false; it would
fail at the Bullet Cluster. The defensible claim distinguishes: dark
\emph{energy} — reading artefact; dark \emph{matter} — genuine
explanatory burden, geometric candidates computed, forward proof open.

% ============================================================
\chapter{Symmetry Check: Does the Argument Cut FFGFT Too?}
% ============================================================

\textbf{[K]} Doc.~267 requires symmetric application. The objection
would be natural: ``FFGFT's fractal time unfolding is just as
computationally present and without independent consequence as
$\Lambda$.'' Checking against the conditions of Ch.~1:

\begin{itemize}
  \item Condition~2 is \emph{violated} for $\xi$ (and the artefact
        status is therefore excluded): $\xi$ is cross-anchored.
        The same ratio carries the lepton mass ladder, the Koide scalar,
        the $\alpha$ reference, $K_\text{frak}$, and the bulk exponent
        ($k^* = \ln\varphi/\xi$, Doc.~275) — and only \emph{thereafter}
        $H_0$ as $\xi^{10}$ power (Doc.~279). A quantity with obligations
        in five foreign sectors cannot become objectless under a
        cosmological reading change; it continues to work elsewhere.
  \item $\Lambda$ has exactly zero such external obligations.
\end{itemize}

This is the decisive asymmetry, and it is bookkeeping, not a value
judgement: not ``our quantity is true'' but ``our quantity has external
obligations, $\Lambda$ has none.''

\medskip
\noindent\textbf{[POSIT] Honesty note:} The magnitude of $\xi$
($5^4$ scale) is declared intuitively/empirically motivated, not a
forward derivation (P33, Doc.~190). The cross-anchoring stands
independently — it concerns the \emph{consequences} of $\xi$, not
its origin.

% ============================================================
\chapter{Demarcation: What This Document Does Not Claim}
% ============================================================

\begin{itemize}
  \item It does not claim that the static reading is proved.
        Doc.~267 stands: no reading is currently empirically preferred.
  \item It does not claim that $\Lambda$CDM computes incorrectly.
        Reading artefacts compute correctly — that is part of the
        definition.
  \item It does not claim that dark matter is fully explained.
        Stages~A–C are computed and passed at order-of-magnitude level.
        Bullet: peak 10~kpc from galaxies (mass-contrast effect).
        Rotation curves: DDO~154 ratio 0.998; NGC~3198/MW ratio $\sim$0.80.
        Open posit: interpolation form $\mu(x)$ from T$^4$ geometry
        not yet derived (P44-B).
  \item It provides no killer test. The contribution is ontological
        sorting: which quantities survive a reading change, which do not
        — and what that says about their status.
\end{itemize}

% ============================================================
\chapter{P44 as Programme: Stage Structure and Results}
% ============================================================

\textbf{[POSIT]} P44 is not a single computation but a three-stage
programme. All three stages have been computed; their status is kept
separate. Five verification scripts, all inputs declared, no fits.

% ----------------------------------------------------------
\section{Stage A — The $a_0$ Anchor}
% ----------------------------------------------------------

\textbf{[K]} The empirical acceleration scale below which rotation
curves flatten is $a_0 \approx 1.2\times10^{-10}$~m/s$^2$
(RAR fit, McGaugh et al.\ 2016 — declared external reference).
From the FFGFT chain Doc.~279,
$H_0/c = (\pi/2)\,\xi^{10}/\bar\lambda_e$, without any further input:
\[
  H_0 = 66.82~\text{km/s/Mpc}, \qquad
  \frac{cH_0}{2\pi} = 1.033\times10^{-10}~\text{m/s}^2,
  \qquad \frac{cH_0/2\pi}{a_0^{\text{emp}}} = 0.86.
\]
The well-known \emph{MOND coincidence} $a_0 \sim cH_0$ is thus in FFGFT
not a coincidence but a consequence of the $\xi^{10}$ chain
(K6, Doc.~190): two independent determination paths
(rotation curve fit; lepton sector via $\bar\lambda_e$ and $\xi$)
meet to within 14\%. Structure analogous to the pion anchor (A155/A270).

\medskip
\textbf{[K] Rotation curves fully computed
(\texttt{ffgft\_307\_p36\_stufeC2\_rotationskurven.py}):}
Correct MOND equation $a_\text{tot}\cdot\mu(a_\text{tot}/a_0) =
a_\text{Newton}$ with $\mu(x) = x/\sqrt{1+x^2}$, $a_0$ from
$\xi^{10}$ chain (no free parameter). Results:

\begin{center}
\begin{tabular}{lccc}
\toprule
Galaxy & $v_\text{Newton}$ & $v_\text{FFGFT}$ & $v_\text{obs}$ / ratio \\
\midrule
DDO~154 (gas-dom.)  & 18.5~km/s & 46.9~km/s & 47.0~km/s / \textbf{0.998} \\
NGC~3198            & 50.4~km/s & 120~km/s  & 150~km/s / 0.80 \\
Milky Way           & 122~km/s  & 175~km/s  & 220~km/s / 0.80 \\
\bottomrule
\end{tabular}
\end{center}

DDO~154 (cleanest test, $M_b$ best determined): ratio 0.998.
NGC~3198 and MW: ratio $\sim$0.80, dominated by $M_b$ uncertainty
($\sim$30\%).

\medskip
\textbf{[K] Status of the interpolation form $\mu$:}
The following is derived from FFGFT:
(i)~$a_0 = cH_0/(2\pi)$ as transition scale ($\xi^{10}$ chain, Stage~A);
(ii)~both limits: $\mu \to 1$ for $a \gg a_0$ (Newton) and
$\mu \to x$ for $a \ll a_0$ (MOND limit $v^4 = GM_b a_0$).
The following is \emph{not} derived: the form of $\mu$ between the limits.
The choice $\mu(x) = x/\sqrt{1+x^2}$ (Simple) is the most natural
analytically-monotone interpolation with the correct limits; many forms
(Standard, Simple, etc.) are empirically indistinguishable without
precise profile measurements.

\medskip
\textbf{[POSIT]} Two honesty notes.
First: the form of $\mu(x)$ between the limits is a posit —
the limits themselves are justified, the interpolation is not
(P44-B open).
Second: the exactly computed chain yields 66.82~km/s/Mpc;
Doc.~280 gives 66.2 — discrepancy 0.9\%, cause to be clarified.

% ----------------------------------------------------------
\section{Stage B — FE Coupling and Selection Computation}
% ----------------------------------------------------------

\textbf{[K]} The mechanism of path elongation is implemented in
Doc.~026 as a finite-element model with ray tracing — it is
\emph{path-dependent by construction} (the ray accumulates delay
element by element along its trajectory).
The prescription for how local mass deforms the element geometry is
\emph{not} an open posit — it is the duality itself. $\tilde T \cdot m = 1$
means: local mass sets the local rate of each element; the intrinsic
time field $T(x)$ is the deformation prescription of the mesh. Sign
automatically correct: element near mass $\to$ $\tilde T$ smaller
$\to$ slower rate $\to$ ray delayed (Shapiro anchor satisfied by
construction).

\medskip
\textbf{[K] The Factor-2 Gate — computed
(\texttt{ffgft\_307\_p36\_stufeB1\_faktor2\_fe.py}):}
FE ray tracing through a radial element grid (30001 log nodes,
linear interpolation, RK4 ray equation). Three configurations of
element deformation $n = 1 + \kappa\,GM/(rc^2)$:

\begin{center}
\begin{tabular}{lcc}
\toprule
Configuration & $\alpha$ at solar limb & Ratio to $1.75''$\\
\midrule
A: time rate only (duality)         & $0.8750''$ & $0.500$\\
B: edge stretching only (path el.)  & $0.8750''$ & $0.500$\\
C: rate $+$ stretching              & $1.7500''$ & $1.000$\\
\bottomrule
\end{tabular}
\end{center}

\textbf{Result:} Duality \emph{alone} fails — it yields the 1911 half
result. Fractal path elongation \emph{alone} yields exactly half too.
Together they yield exactly the measured value.

\medskip
\textbf{[S] Two-Halves Synthesis (candidate):} The two mechanisms of
the corpus are the two halves of the Einstein deflection —
$\tilde T \cdot m = 1$ the time half (1911), fractal path elongation
the space half (1911 $\to$ 1915). They do not replace each other;
they are complementary.

\medskip
\textbf{[K] Selection computation under T$^4$ invariance rules
(\texttt{ffgft\_307\_p36\_stufeB2\_kappa\_herleitung.py}):}
Starting point $m(r) = m_0(1+\Phi/c^2)$ (dual reading of the measured
gravitational time dilation, Pound--Rebka anchored; mass decreases in
the potential well), rate $\tilde T = 1/m$, so $\kappa_\text{time} = 1$
fixed. Four candidates for the spatial edge rule, all tested in the same
FE setup:

\begin{center}
\begin{tabular}{lccc}
\toprule
Rule & $\gamma$ & $\alpha$ & Verdict\\
\midrule
K1 conformal ($l = \lambda_C$, same source) & $-1$ & $0.0000''$ & excl.\\
K2 unimodular (4-vol.\ invariant)           & $+1/3$ & $1.1667''$ & excl.\\
K3 area (time-space 2-cell invariant)       & $+1$ & $1.7500''$ & \textbf{passes}\\
K0 rate only (1911 reference)               & $0$ & $0.8750''$ & excl.\\
\bottomrule
\end{tabular}
\end{center}

\textbf{The conformal trap:} The most natural path — edge length equal
to Compton length $\hbar/mc$ of the \emph{same} node — makes the geometry
conformally flat: light sees nothing, $\alpha = 0$. Harder excluded than
the half result.

\medskip
\textbf{[S] The survivor and its form:} K3 reads
$L_t \cdot l = \text{const}$ — area invariance of the time-space 2-cell.
This is formally the reciprocity form of $\tilde T \cdot m = 1$, applied
a second time, now to the time-space pair; the corpus formula ``$c$ =
space-time bridge'' thereby receives a precise reading: $c$ as the
invariant area of the cell.
Consequence: $\gamma_\text{PPN} = 1$ \emph{exactly} (Cassini passed with
unlimited margin; K1/K2 lie a factor $10^5$ resp.\ $3\times10^4$ above the
bound). Shapiro consistency numerically validated (rel.\ dev.\ $8\times10^{-5}$
against the analytic path-integral form).

\medskip
\textbf{[POSIT] Status and limits:} This is an exclusion proof over a
\emph{declared} candidate list plus one exactly passing survivor — not
a forward proof. Remaining open: why the 2-cell area (and not the
4-volume) is the invariant of the T$^4$ metric. Completeness of the
candidate list is a posit.

\medskip
\textbf{[POSIT] Safeguards against self-deception:}
(i) The characteristic scale $r_\xi(M)$ must \emph{not} be set to the
Compton wavelength of the total mass ($\hbar/Mc \sim 10^{-68}$~m for a
galaxy) — with rescue exponents that would be free-parameter renaming;
the scale must follow from the Doc.~182 geometry. (ii) Expressions of
the form $\Phi_\text{frak} \propto \int (d\xi/dr)\,dr$ are empty
identities without derivational content and will not be admitted to the
corpus. Profile candidates without derivation are, if used, to be
declared as posits with their own parameter.

% ----------------------------------------------------------
\section{Stage C — Bullet Cluster and Rotation Curves: Computed}
% ----------------------------------------------------------

\textbf{[K]} Bullet Cluster requirement from published lensing maps
(Clowe et al.\ 2006, ApJ 648 L109): lensing mass at the subcluster
galaxy position $\sim 2.3\times 10^{14}\,M_\odot$, baryons there
$\sim$10\%, spatial offset lensing peak -- gas peak $\sim$250~kpc.

\medskip
\textbf{[POSIT]} For Doc.~182 the expectation was: the $\xi$-field
profile would follow the gas density at instantaneous coupling —
and thus fail at the Bullet. The computation shows the opposite.

\medskip
\textbf{[K] Sharpening:} Pre-collision, the baryonic potential of both
clusters is \emph{gas-dominated} ($\sim$85\% of baryons are ICM gas;
stars/galaxies $\sim$15\%). A pure potential-\emph{memory} field would
remember gas-dominated potential wells — it would smear at the collision
site, not peak at the \emph{current} galaxy position. Memory alone is
insufficient: the field would need to be \emph{advected} with the
collision-free component.

\medskip
\textbf{[K] Test setup and result
(\texttt{ffgft\_307\_p36\_stufeC1\_bullet\_cluster.py}):}
Two-clump mesh, projection along line of sight (weak lensing),
$\tilde T \cdot m = 1$ (K3, instantaneous coupling). Parameters
directly from Clowe+2006: $M_\text{gal} = 2.3\times10^{14}\,M_\odot$,
$M_\text{gas} = 0.23\times10^{14}\,M_\odot$, offset 250~kpc.

\begin{center}
\begin{tabular}{lcc}
\toprule
Quantity & Value & Requirement \\
\midrule
Lensing peak position (FFGFT) & $x = 10.0$~kpc & at galaxies \\
Distance peak -- galaxies & 10~kpc & $< 50$~kpc \\
Distance peak -- gas & 240~kpc & $> 150$~kpc \\
\bottomrule
\end{tabular}
\end{center}

\textbf{[K] Verdict: Bullet test passed} — unexpected at instantaneous
coupling. \textbf{Cause:} $M_\text{gal}/M_\text{gas} = 10$ (observed).
The galaxy potential dominates the gas field despite a 250~kpc offset.
Sensitivity test with $M_\text{gal} = M_\text{gas}$ (pessimistic):
peak at 30~kpc — still at the galaxies. The profile radius decides,
not the advection history.

\medskip
\textbf{[POSIT] Honesty note:} This is a 1D projection test, not a
full 2D lensing reconstruction. The decisive factor is the mass ratio 10,
not the mechanism itself. A sharper test requires: (i) observed 2D lensing
profile against computation; (ii) test with sub-clusters where
$M_\text{gal}/M_\text{gas}$ is smaller. The collapse risk of field-DM
is absent at instantaneous coupling: the field follows the instantaneous
density, no advected field needed.

% ============================================================
\chapter{Layer-Status Summary}
% ============================================================

\begin{itemize}
  \item \textbf{[K]} Degeneracy of the three data pillars: taken from
        Doc.~267.
  \item \textbf{[K]} $\Lambda$ satisfies all three artefact conditions
        (Ch.~2).
  \item \textbf{[K]} DM violates Condition~1; evidence reading-independent
        (Ch.~3).
  \item \textbf{[K]} Asymmetry $\xi$ vs.\ $\Lambda$ via cross-anchoring
        (Ch.~4).
  \item \textbf{[K]} Stage A: $a_0$ anchor $cH_0/2\pi =
        1.033\times10^{-10}$~m/s$^2$, 14\% deviation from RAR value
        (Ch.~6, §1).
  \item \textbf{[K]} Stage B: factor-2 gate (each $0.875''$ time and
        space, together $1.750''$); K3 sole survivor of the selection
        computation, $\gamma_\text{PPN}=1$ exactly (Ch.~6, §2).
  \item \textbf{[K]} Stage C: Bullet test passed (peak 10~kpc from
        galaxies, mass-contrast effect). Rotation curves: DDO~154 ratio
        0.998; NGC~3198/MW ratio $\sim$0.80 ($M_b$ uncertainty dominates).
        $\mu$ limits justified, interpolation form posit (Ch.~6, §3).
  \item \textbf{[POSIT]} Artefact definition (Ch.~1); DM as open
        geometric problem with candidate Doc.~182 (Ch.~3); P44 programme
        A/B/C fully computed; $\mu$ limits justified, interpolation form
        and K3 forward proof open (Ch.~6).
\end{itemize}

% ============================================================
\chapter{Open Points (Doc.~190 P44)}
% ============================================================

\begin{enumerate}
  \item \textbf{P44-B remaining:} Factor-2 gate computed (time$+$space
        each $0.875''$, together $1.750''$); selection computation
        completed (K3 area invariance as sole survivor, $\gamma = 1$
        exactly). Remaining: forward proof that the time-space 2-cell
        area (and not the 4-volume) is the invariant of the T$^4$ metric
        — plus completeness check of the candidate list.
  \item \textbf{P44-C Bullet:} Computed
        (\texttt{ffgft\_307\_p36\_stufeC1\_bullet\_cluster.py}).
        Contrary to expectation, passed: $M_\text{gal}/M_\text{gas}=10$
        dominates at instantaneous coupling. Lensing peak 10~kpc from
        galaxies, 240~kpc from gas. Collapse risk field-DM absent
        (instantaneous coupling without advection). Open refinement:
        2D lensing reconstruction, sub-clusters with smaller mass contrast.
  \item \textbf{P44-C Rotation curves:} Computed
        (\texttt{ffgft\_307\_p36\_stufeC2\_rotationskurven.py}).
        DDO~154: ratio 0.998 (passed). NGC~3198/MW: ratio $\sim$0.80
        ($M_b$ uncertainty dominates, not mechanism). Open posit: form
        of $\mu(x)$ between the limits not derived from T$^4$ geometry
        (limits themselves justified, interpolation open — P44-B).
  \item \textbf{$H_0$ discrepancy:} 66.82 vs.\ 66.2~km/s/Mpc —
        cause to be clarified in Doc.~279/280 (rounding/$m_e$ version).
  \item Relation of the artefact analysis to the early universe phase
        (BBN, structure formation, Doc.~267 §\,honest position) — there,
        one reading might yet acquire reading-independent burdens.
  \item Check whether the artefact definition (Ch.~1) applies to further
        $\Lambda$CDM quantities (inflaton parameters?) — and symmetrically
        to FFGFT quantities in the reduction layer.
\end{enumerate}

% ============================================================
\chapter{References and External Sources}
% ============================================================

\textbf{Internal documents:}
Doc.~267 (cosmological degeneracy);
Doc.~280 (achromatic static redshift, K6);
Doc.~182 (fractal path elongation, Doc.~026 FE implementation);
Doc.~279 ($H_0$ as $\xi$ power, $\Lambda$ fine-tuning in outlook);
Doc.~275 ($k^*$ recursion);
Doc.~190 (P18--P20, P33, P35, P44, K6).

\medskip
\noindent\textbf{Verification scripts:}
\texttt{ffgft\_307\_p36\_stufeA\_a0\_anker.py} (all inputs declared, no fits);
\texttt{ffgft\_307\_p36\_stufeB1\_faktor2\_fe.py} (30001 log nodes, RK4, convergence checked);
\texttt{ffgft\_307\_p36\_stufeB2\_kappa\_herleitung.py} (selection computation, four candidates);
\texttt{ffgft\_307\_p36\_stufeC1\_bullet\_cluster.py} (two-clump projection, Clowe+2006 parameters);
\texttt{ffgft\_307\_p36\_stufeC2\_rotationskurven.py} (DDO~154, NGC~3198, MW, no free parameters).

\medskip
\noindent\textbf{Declared external references:}
McGaugh, Lelli, Schombert 2016, PRL 117, 201101 (RAR/$a_0$);
Clowe et al.\ 2006, ApJ 648, L109 (Bullet Cluster).
